\chapter{Повторы: циклы}\label{ch:cikly}

\textbf{Что вы узнете:} как не писать одно и то же много раз.

\section{Цикл в программе одного дрона}

Три красные вспышки можно записать шестью строками, а можно одним циклом:

\solfile{l13_blink.py}

\kod{for _ in range(3):} значит «повтори три раза». Всё, что в следующих строках записано \emph{со сдвигом вправо} (четыре пробела), — тело цикла, оно повторяется. Сдвиг вправо в Python обязателен: так он понимает, где кончается цикл.

Подчёркивание \kod{_} здесь вместо имени счётчика: нам не важен номер повтора. Если номер нужен, берут имя: \kod{for n in range(3):} даст \kod{n} = 0, 1, 2.

\section{Цикл по списку}

\kod{for color in colors:} перебирает значения списка по одному:

\begin{verbatim}
colors = ['red', 'green', 'blue']
for color in colors:
    d.rgb(color)
    d.set_z(100, 1)
\end{verbatim}

\section{Цикл по дронам}

Если много дронов делают похожее, цикл можно поставить \emph{вокруг} \kod{show.drone}:

\solfile{t10_loop.py}

Каждый проход создаёт очередного дрона: \kod{show.drone(n + 1)}. Список \kod{corners} хранит клетки, \kod{colors} цвета. Выражение \texttt{corners[(n + k) \% 4]} берёт клетку по номеру, а \texttt{\%} (остаток от деления) возвращает счёт к нулю после четвёртой клетки: после номера 3 идёт 0.

\begin{vnimanie}[Циклы и программа дрона]
Обычно вкладка «Код» показывает программу одного дрона. Если же дроны создаёт общий цикл, отдельной программы у каждого нет, и приложение показывает код целиком: «Дроны этого шоу создаёт общий цикл, поэтому код показан целиком». Это не ошибка: симуляция, партитура и проверки работают как обычно. Цикл \emph{внутри} программы одного дрона такому показу не мешает.
\end{vnimanie}

\begin{probuy}
  \item «Челнок»: три раза туда и обратно между клетками 13 и 17 (по 4\,с), циклом.
  \item Радуга: список из семи цветов \kod{['red', 'yellow', 'green', 'cyan', 'blue', 'purple', 'white']} и цикл по нему, по секунде каждый.
  \item Три дрона в клетках 13, 15, 17 мигают красным по очереди (по 2 секунды). Пауза зависит от номера дрона. Запишите циклом по дронам.
  \item Квадрат три раза подряд: цикл \kod{for lap in range(3):} с \kod{code_s}.
\end{probuy}
\otvety{13}
